\documentclass[12pt,letterpaper]{article}
\usepackage[T1]{fontenc}
\usepackage{lmodern,amsmath,amssymb,geometry,xcolor,hyperref,fancyhdr,booktabs}
\geometry{left=.92in,right=.92in,top=.9in,bottom=.85in,headheight=15pt,headsep=16pt}
\definecolor{ink}{HTML}{1B3340}
\definecolor{accent}{HTML}{17665E}
\definecolor{muted}{HTML}{4A6068}
\hypersetup{colorlinks=true,linkcolor=accent,urlcolor=accent,
 pdftitle={Compositions and a divisibility theorem for OEIS A115047},
 pdfauthor={Research notes for Blake Stone, developed with OpenAI Codex},
 pdfsubject={Proofs concerning OEIS A115047, A208354, and A323111}}
\pagestyle{fancy}\fancyhf{}
\lhead{\small\sffamily\color{muted}Compositions and divisibility}
\rhead{\small\sffamily\color{accent}OEIS research notes}
\cfoot{\small\sffamily\color{muted}\thepage}
\renewcommand{\headrulewidth}{.25pt}
\setlength{\parindent}{0pt}\setlength{\parskip}{6pt}
\setlength{\abovedisplayskip}{7pt}\setlength{\belowdisplayskip}{7pt}
\setlength{\abovedisplayshortskip}{4pt}\setlength{\belowdisplayshortskip}{5pt}
\linespread{1.07}
\newcommand{\tagline}[1]{{\small\sffamily\bfseries\color{accent}#1}\par}
\newcommand{\titleline}[1]{{\fontsize{23}{27}\selectfont\sffamily\bfseries\color{ink}#1}\par}
\newcommand{\subhead}[1]{\vspace{1pt}{\large\sffamily\bfseries\color{ink}#1}\par}
\newcommand{\lead}[1]{\textbf{\color{ink}#1.}\enspace}
\newcommand{\qedx}{\hfill\(\square\)}
\newcommand{\source}[1]{\vfill{\small\color{muted}#1\par}}
\newcommand{\oeis}[1]{\href{https://oeis.org/#1}{OEIS #1}}
\newcommand{\ph}{\varphi}
\begin{document}
\thispagestyle{plain}
\tagline{A115047 \enspace / \enspace A208354 \enspace / \enspace A323111}
\titleline{Compositions and a divisibility theorem\\for OEIS A115047}
{\small\sffamily Research notes for Blake Stone \enspace$\cdot$\enspace Developed with OpenAI Codex
 \hfill November 26, 2026}\par
\vspace{6pt}
\subhead{1\quad The smallest universal multiplier}
Let $v_n$ be A115047, indexed from $n=3$. Michael Somos proposed in 2014 that
every prime index divides its term. The first term is $v_3=1$. The following
statement identifies the sole prime exception and extends to composite indices.

\lead{Theorem 1} For every $n\ge3$,
\[
                         n\mid 12v_n .
\]
Equivalently, $n/\gcd(n,12)$ divides $v_n$. The positive integers $C$ satisfying
$n\mid Cv_n$ for every $n\ge3$ are exactly the multiples of $12$. In particular,
$p\mid v_p$ for every prime $p\ge5$.

\lead{Proof} Zograf's recurrence, with $j=i+1$, is
\[
 2(n-1)v_n=\sum_{j=2}^{n-2}\binom{n-4}{j-2}\binom{n}{j}
 (j-1)(n-j-1)v_{j+1}v_{n-j+1}\qquad(n\ge4).
\]
The $v_n$ are integers by the border-strip enumeration of Alexandersson and
Jordan. Put $D_n=(n-1)(n-2)(n-3)$. Cancellation of factorials gives
\[
 D_n\binom{n-4}{j-2}\binom{n}{j}
 =n(j-1)(n-j-1)\binom{n-1}{j-1}\binom{n-1}{j}.
\]
After multiplying the recurrence by $D_n$, every summand on the right is
divisible by $n$. On the left, the multiplier of $v_n$ satisfies
\[
 2(n-1)D_n=2(n-1)^2(n-2)(n-3)\equiv 2(-1)^2(-2)(-3)=12\pmod n .
\]
Hence $n\mid 12v_n$ for $n\ge4$, and the statement holds at $n=3$ as well.
The initial value and the recurrence give $v_3=v_4=1$. Thus any universal $C$
must be divisible by both $3$ and $4$, which proves optimality. For every
prime $p\ge5$, $\gcd(p,12)=1$ gives the asserted prime divisibility.\qedx

\source{\oeis{A115047}. Michael Somos's conjecture is dated September 20, 2014.
For the recurrence and integrality, see P.~Alexandersson and L.~Jordan,
\href{https://cs.uwaterloo.ca/journals/JIS/VOL22/Alexandersson/alex4.pdf}
{\textit{Enumeration of Border-Strip Decompositions and Weil-Petersson Volumes}},
J.\ Integer Sequences 22 (2019), Article 19.4.5, Theorems 32 and 34.}

\newpage
\tagline{A208354 \enspace / \enspace COMPOSITIONS}
\titleline{2\quad One even part or one run of ones}
A composition of $n$ is an ordered list of positive integers summing to $n$.
There is one empty composition of zero.

\lead{Theorem 2} For each $n\ge0$, as many compositions of $n$ have at most one
even part as have all their ones in a single consecutive run. The latter
class includes compositions with no ones.

\lead{Proof} Both classes have a weight-preserving bijection to words in
$1,2,3$ of total weight $n$ with at most one $3$. This common class is
Jerrold Grossman's interpretation of A208354. The same construction works in
both directions after exchanging the roles of $1$ and $2$.

Take $(a,b)=(1,2)$ or $(2,1)$. An object is a list of ordinary blocks of weight
$a+kb$ ($k\ge0$), with at most one distinguished block of weight $tb$ ($t\ge1$).
For $(1,2)$ the distinguished block is the even part. For $(2,1)$ it is
the whole run of $t$ ones, while the ordinary blocks are the parts at least two.
Let $U,V,*$ be letters with respective weights $a,b,a+b$.

\lead{Encoding} If the distinguished block is absent or last, write $V^kU$
for each ordinary block and append $V^t$ if the distinguished block is present.
Otherwise, encode all ordinary blocks before the distinguished block as $V^kU$.
Encode the distinguished block and the ordinary block immediately after it by
\[
             (tb,\ a+kb)\ \longmapsto\ V^{t-1}\,*\,V^k .
\]
Encode each later ordinary block of weight $a+hb$ as $UV^h$. The displayed
replacement preserves weight since $*$ contributes $a+b$.

\lead{Decoding} In a word without $*$, each $U$ ends an ordinary block.
Trailing $V$'s, if present, give the terminal distinguished block. In a word
with $*$, parse the left side up through its last $U$, treating each $U$ as a
block end. Its remaining $V$'s give $t-1$. The initial $V$'s on the right
give $k$. Each remaining $U$ starts a later ordinary block. This decoding
also works when either side is empty. It recovers exactly one object, so the two
encodings give the claimed bijection.\qedx

For example, the composition $(3,4,1,3)$ in the first class and $(2,3,1,3,2)$
in the second both encode as $(2,1,2,3,1,2)$. Each has total weight $11$.
The common generating function follows by splitting a word at its optional $3$.
With $D=(1-x-x^2)^{-1}$ it is
\[
 D+x^3D^2=\frac{(1-x)(1-x^2)}{(1-x-x^2)^2}.
\]
\source{\oeis{A208354}. Gregory L.~Simay proposed the single-run interpretation
on September 29, 2018. Jerrold Grossman recorded the interpretation using a
single optional $3$ on April 3, 2024.}

\newpage
\tagline{A323111 \enspace / \enspace ASYMPTOTIC ENUMERATION}
\titleline{3\quad Compositions with range $n$}
Let $a_n$ count compositions of $2n$ whose largest part exceeds their smallest
part by $n$. Write $\ph=(1+\sqrt5)/2$.

\lead{Theorem 3} As $n\to\infty$,
\[
 a_n=(n+2)2^{n-3}-\frac{n\ph^n}{5\ph^5}+O(\ph^n).
\]
In particular $a_n\sim n2^{n-3}$, as conjectured by V\'{a}clav Kot\v{e}\v{s}ovec in 2019.

\subhead{Removing the unique largest part}
If the smallest part is $m$, the largest part $n+m$ exceeds half the total
weight. It occurs exactly once. Removing it leaves two possibly empty lists,
one on each side, with combined weight $n-m$. Every remaining part is at
least $m$, and at least one must equal $m$. Conversely, any such pair gives a
valid composition upon reinsertion of $n+m$ between its lists. No remaining
part can exceed the total remaining weight $n-m<n+m$.

Let $b_m(s)$ count pairs of ordered lists of combined weight $s$ with parts at
least $m$. Set $b_m(0)=1$ and $b_m(s)=0$ for $s<0$. The bijection gives
\[
 a_n=\sum_{m\ge1}\bigl(b_m(n-m)-b_{m+1}(n-m)\bigr).                 \tag{1}
\]
For one list the allowed part weights have series $x^m/(1-x)$. Summing over
list lengths and then taking an ordered pair gives
\[
 B_m(x)=\sum_{s\ge0}b_m(s)x^s=\left(1-\frac{x^m}{1-x}\right)^{-2}.    \tag{2}
\]
In particular $xB_1(x)=x(1-x)^2/(1-2x)^2$, whose $x^n$ coefficient equals
$(n+2)2^{n-3}$ for $n\ge2$.

\subhead{Why this gives the leading term}
If $m\ge2$, the number $c_m(s)$ of single lists of weight $s$ is at most
$\ph^s$. Indeed $c_m(0)=1$ and induction on the first part gives
$c_m(s)=\sum_{j=m}^{s}c_m(s-j)\le\ph^s\sum_{j\ge2}\ph^{-j}=\ph^s$
for $s>0$. Splitting the weight between two lists gives
$b_m(s)\le(s+1)\ph^s$ for $s\ge0$. In (1) the absolute error from keeping
only $b_1(n-1)$ is at most
\[
 b_2(n-1)+\sum_{m=2}^{n} b_m(n-m)=O(n\ph^n),
\]
by geometric summation. Since $\ph<2$, this proves the proposed asymptotic.
The next page determines the first exponentially smaller correction.

\source{\oeis{A323111}. Alois P.~Heinz defined the sequence on January 4, 2019.
V\'{a}clav Kot\v{e}\v{s}ovec proposed the leading asymptotic on January 7, 2019.}

\newpage
\tagline{A323111 \enspace / \enspace REFINEMENT}
\titleline{The coefficient of the correction term}
Summing (1) over $n\ge1$ and shifting the index in the second sum gives an
identity of formal power series,
\[
 \sum_{n\ge1}a_nx^n
 =x B_1(x)-x(1-x)B_2(x)-(1-x)\sum_{m\ge3}x^{m-1}B_m(x).             \tag{3}
\]
Each coefficient uses finitely many summands.

\lead{Bounding the tail} Let $\lambda>1$ be the real root of
$\lambda^3=\lambda^2+1$. Then $\sum_{j\ge3}\lambda^{-j}=1$ and
$\lambda<\ph$. The induction on page 3 gives $b_m(s)\le (s+1)\lambda^s$
uniformly for $m\ge3$ and $s\ge0$. The absolute value of the $x^n$
coefficient of the last term of (3) is at most
\[
 \sum_{m=3}^{n+1} \bigl(b_m(n-m+1)+b_m(n-m)\bigr)=O(n\lambda^n).
\]
This estimate is $o(\ph^n)$.

\lead{The second term} The rational function subtracted in the middle of (3)
has a double pole at $\rho=1/\ph$,
\[
 R(x)=x(1-x)B_2(x)
 =\frac{x(1-x)^3}{(1-\ph x)^2(1+x/\ph)^2}.
\]
The coefficient of $(1-\ph x)^{-2}$ in its partial fractions is
\[
 K=\left.(1-\ph x)^2R(x)\right|_{x=\rho}
  =\frac{\rho(1-\rho)^3}{(1+\rho^2)^2}
  =\frac{\rho^5}{5},
\]
using $1-\rho=\rho^2$ and $1+\rho^2=\sqrt5\,\rho$. The simple-pole term at
$\rho$ contributes $O(\ph^n)$ to the $n$th coefficient. The other pole at
$-\ph$ contributes $O(n\ph^{-n})$. Since
$[x^n](1-\ph x)^{-2}=(n+1)\ph^n$, we obtain
$[x^n]R=Kn\ph^n+O(\ph^n)$. Inserting this in (3) proves Theorem 3.\qedx

\subhead{Numerical check}
The exact integer counts in (1) give the following values of
$((n+2)2^{n-3}-a_n)/(n\ph^n)$. The theorem predicts their limit
$K=1/(5\ph^5)$.

\begin{center}
\begin{tabular}{rrrr}
\toprule
$n=100$ & $n=500$ & $n=1000$ & $K$\\
\midrule
0.019009366 & 0.018229004 & 0.018131497 & 0.018033989\\
\bottomrule
\end{tabular}
\end{center}
\source{The companion archive includes the typesetting source and a standard-library
Python verifier. It checks the volume recurrence against the multinomial formula,
both composition bijections against word enumeration, and (1) against direct
enumeration of compositions of $2n$.}
\end{document}
